\chapter{情報理論}

\section{ハミング距離}
ハミング距離は、ある文字列を別の文字列に変形する際に必要な置換回数を数えたものである。
ただし、文字列の長さは等しい。
例をいくつかあげる。
\begin{itemize}
  \item 101\RED{0}1\RED{0}と101\RED{1}1\RED{1}のハミング距離は2である。
  \item 1\RED{234}5と1\RED{111}5のハミング距離は3である。
  \item \RED{ho}g\RED{e}と\RED{fu}g\RED{a}のハミング距離は3である。
\end{itemize}


\section{グレイコード}
グレイコードは数値の符号化法のひとつで、前後に隣接する符号間のハミング距離が必ず1であるという特性を持つ。

以下に、2,3bitのグレイコードを示す。
\begin{table}[htb]

  \caption{グレイコード }
  \label{table:GrayCode}

  \centering
  \begin{tabular}{c|cc}
      & 2bit & 3bit \\ \hline
    0 &  00  & 000 \\
    1 &  01  & 001 \\
    2 &  11  & 011 \\
    3 &  10  & 010 \\
    4 &      & 110 \\
    5 &      & 111 \\
    6 &      & 101 \\
    7 &      & 100
  \end{tabular}

\end{table}

10進数$n$をグレイコードにするためには、「$n$の二進数表現」と「これを1ビット右シフトし、先頭に0をつけたもの」との排他的論理和$\oplus$とる。
例えば、10進数の10をグレイコードにしてみる。
\begin{enumerate}
  \item まず、10の二進数表現は$1010_{(2)}$
  \item これを1ビット右にシフトし、先頭に0をつけると、$0101_{(2)}$
  \item XORをとると、$1010 \oplus 0101 = 1111$
\end{enumerate}

C言語で書くと$n$をint型の変数とすると、$n\text{\textasciicircum}(n>>1)$でグレイコードに変換できる。

グレイコードにも順序があり、それはグレイコードを整数に戻したときの、整数の順序である。
なので、グレイコードでの$11$は整数では$2$で、$10$は$3$なので、$11$の方が$10$より先である。
